%\documentclass[12pt,letterpaper]{article}



%%%
%% Required packages:
%%
\usepackage{etex}
\usepackage{amsmath}
\usepackage{amssymb}
\usepackage{amstext}
\usepackage{amsthm}
\usepackage{fancyhdr,fancybox}
\usepackage{mathrsfs} % need for \mathscr command
\usepackage{multicol}
\usepackage[normalem]{ulem}
%\usepackage{pictex,pictexplus}
% \usepackage{pictexwd}
\usepackage{rotating}
\usepackage{url}
\usepackage{xfrac}
\usepackage{booktabs}
\usepackage{color,soul}
\usepackage{comment}

\usepackage{hyperref}
\hypersetup{
    colorlinks=true,     % false: boxed links; true: colored links
    linkcolor=blue,      % color of internal links
    citecolor=blue,      % color of links to bibliography
    filecolor=blue,      % color of file links
    urlcolor=blue        % color of external links
}


\usepackage{tikz}
\usetikzlibrary{backgrounds,calc}

%%
%% Common formatting commands
%%
\setlength{\parindent}{0in}
\setlength{\textwidth}{7in}
\setlength{\evensidemargin}{-0.25in}
\setlength{\oddsidemargin}{-0.25in}
\setlength{\parskip}{.5\baselineskip}
\setlength{\topmargin}{-0.5in}
\setlength{\textheight}{9in}

%               Problem and Part
\newcounter{problemnumber}
\newcounter{partnumber}
\newcounter{subpartnumber}
\newcommand{\Problem}{%
  \refstepcounter{problemnumber}%
  \setcounter{partnumber}{0}%
  \setcounter{subpartnumber}{0}%
  \item[\makebox{\hfil\textbf{\theproblemnumber.}\hfil}]%
  }
\newcommand{\InProblem}[2][1.6]{%
  \refstepcounter{problemnumber}%
  \setcounter{partnumber}{0}%
  \setcounter{subpartnumber}{0}%
  \textbf{\theproblemnumber.}\ \ \parbox[t]{#1in}{#2}%
}
\newcommand{\InSmallProblem}[1]{\InProblem[1.05]{#1}}
\newcommand{\NoProblem}[1]{%
  \mbox{\hfil\textbf{#1}\hfil}%
  }
\newcommand{\UnProblem}{%
  \refstepcounter{problemnumber}%
  \setcounter{partnumber}{0}%
  \setcounter{subpartnumber}{0}%
  \mbox{\hfil\textbf{\theproblemnumber.}\hfil}%
  }
\newcommand{\InFlexProblem}[1]{%
  \refstepcounter{problemnumber}%
  \setcounter{partnumber}{0}%
  \setcounter{subpartnumber}{0}%
  \textbf{\theproblemnumber.}\ \ #1%
}
%%  Part:
\newcommand{\Part}{%
  \renewcommand{\thepartnumber}{(\alph{partnumber})}%
  \refstepcounter{partnumber}
  \setcounter{subpartnumber}{0}%
  \item[(\alph{partnumber})]%
}
\newcommand{\InPart}[2][1.75]{%
  \renewcommand{\thepartnumber}{(\alph{partnumber})}%
  \refstepcounter{partnumber}%
  \setcounter{subpartnumber}{0}%
  (\alph{partnumber})\ \ \parbox[t]{#1in}{#2}%
}
\newcommand{\UnPart}{%
  \refstepcounter{partnumber}%
  \renewcommand{\thepartnumber}{(\alph{partnumber})}%
  \setcounter{subpartnumber}{0}%
  (\alph{partnumber})%
}
\newcommand{\InLargePart}[1]{\InPart[3.05]{#1}}
\newcommand{\InFullPart}[1]{\InPart[6.2]{#1}}
\newcommand{\InSmallPart}[1]{\InPart[1.05]{#1}}
%% SubPart:
\newcommand{\SubPart}{%
  \refstepcounter{subpartnumber}%
  \item[(\roman{subpartnumber})]%
  }
\newcommand{\InSubPart}[2][2.25]{%
  \stepcounter{subpartnumber}%
  (\roman{subpartnumber})\ \ \parbox[t]{#1in}{#2}%
}
\newcommand{\InSmallSubPart}[1]{\InSubPart[1.5]{#1}}
\newcommand{\InTinySubPart}[1]{%
  \stepcounter{subpartnumber}%
  (\roman{subpartnumber})\ \ \makebox{#1}%
}
\newcommand{\InMySubPart}[2][2.25]{%
  \stepcounter{subpartnumber}%
  \parbox[t]{#1in}{\makebox[0.3in][l]{(\roman{subpartnumber})}\ \ {#2}}%
}
\newcommand{\TrueFalse}{\textbf{T}\ \ \textbf{F}\ \ }
\newcommand{\TFPart}[1]{% 
  \stepcounter{partnumber}%
  \setcounter{subpartnumber}{0}%
  \item[(\alph{partnumber})] \TrueFalse\parbox[t]{5.5in}{#1}}

\newcommand{\Points}[1]{(#1 points) \ }
\newcommand{\Point}[1]{(#1 point) \ }


%% For Answers:
\newcounter{answernumber}
\newcommand{\Answer}{%
  \refstepcounter{answernumber}%
  \setcounter{partnumber}{0}%
  \setcounter{subpartnumber}{0}%
  \item[\makebox{\hfil\textbf{\theanswernumber.}\hfil}]%
  }


% Setting up the headers for the first page
%\chead{} % will default to what is typed in the file.
\fancyhead[L]{\textsc{Math 22a}}
\fancyhead[R]{\textsc{Fall 2024}}
\fancyfoot[R]{
	\vspace*{\fill}\footnotesize\textcopyright{} \emph{Harvard Math 22a}	
}
\renewcommand{\headrulewidth}{0pt}

%\fancyhead[C]{}
%	\chead{}	
%	\lhead{}
%	\rhead{}
%	\cfoot{}


% Putting copyright on the footer of each page
%\fancypagestyle{secondpage}{
%	\fancyhead[C]{TESTing again} % change this to be blank
%		%\chead{}	
%	\fancyhead[L]{bears} % change this to be blank
%	\fancyhead[R]{dogs} % change this to be blank
%	\fancyfoot[R]{ %keeps the footer the same
%		\vspace*{\fill}\footnotesize\textcopyright{} \emph{Harvard Math 22a}	
%	}
%}
%\renewcommand{\headrulewidth}{0pt}
%\pagestyle{fancy}
\pagestyle{empty}


%\lhead{}
%\rhead{}
%\rfoot{\vspace*{\fill}\footnotesize\textcopyright{} \emph{Harvard Math 22a}}
%\renewcommand{\headrulewidth}{0pt}
%\pagestyle{fancy}




%%
%% Common shortcut commands:
%%
\newcommand{\ds}{\displaystyle}
\newcommand{\sgn}{\operatorname{sgn}}
\renewcommand{\Pr}{\operatorname{P}}
\newcommand{\CPr}[3][{}]{\Pr_{\text{#1}}\left(#2\,|\,#3\right)}

\newcommand{\C}{\mathbb{C}}
\newcommand{\R}{\mathbb{R}}
\newcommand{\Z}{\mathbb{Z}}
\newcommand{\N}{\mathbb{N}}
\newcommand{\Q}{\mathbb{Q}}
\newcommand{\D}{\mathbb{D}}

\newcommand{\rref}{\operatorname{rref}}
\newcommand{\im}{\operatorname{im}}
\newcommand{\rank}{\operatorname{rank}}
\newcommand{\diag}{\operatorname{diag}}
\newcommand{\Span}{\operatorname{span}}
%\renewcommand{\vec}[1]{\mathbf{#1}}
\newcommand{\charpoly}{\mathscr{P}}
\newcommand{\nicefrac}[2]{\sfrac{#1}{#2}} % using xfrac package
\newcommand{\topology}{\mathcal{T}}
\newcommand{\Inn}{\operatorname{Inn}}

\newcommand{\blankbox}[2]{\fbox{\rule{#1}{0in}\rule{0in}{#2}}}

\newenvironment{myindentpar}[1]%
 {\begin{list}{}%
         {\setlength{\leftmargin}{#1}
          \setlength{\rightmargin}{\leftmargin}}%
         \item[]%
 }
 {\end{list}}
\newtheorem{theorem}{Theorem}
\newtheorem*{NStheorem}{Theorem}
\newtheorem{cor}[theorem]{Corollary}
\newtheorem*{claim}{Claim}
\newtheorem{defn}{Definition}

\newenvironment{amatrix}[1]{%
  \left(\begin{array}{@{}*{#1}{c}|c@{}}
}{%
  \end{array}\right)
}

%--------grstep
% For denoting a Gauss' reduction step.
% Use as: \grstep{\rho_1+\rho_3} or \grstep[2\rho_5 \\ 3\rho_6]{\rho_1+\rho_3}
\newcommand{\grstep}[2][\relax]{%
   \ensuremath{\mathrel{
       {\mathop{\longrightarrow}\limits^{#2\mathstrut}_{
                                     \begin{subarray}{l} #1 \end{subarray}}}}}}
\newcommand{\swap}{\leftrightarrow}


%\usetikzlibrary{arrows}
%\usepackage{wrapfig}
%
%\makeatletter
%\renewcommand*\env@matrix[1][*\c@MaxMatrixCols c]{%
	%\hskip -\arraycolsep
	%\let\@ifnextchar\new@ifnextchar
	%\array{#1}}
%\makeatother
%
%\def\env@matrix{\hskip -\arraycolsep
	%\let\@ifnextchar\new@ifnextchar
	%\array{*\c@MaxMatrixCols c}}

\chead{\Large\textbf{Eigenvalues and Eigenvectors, $\vec\S$ 5.4}}% 

%\begin{document}
\thispagestyle{fancy}


Recall from last class:

\begin{itemize}
\item {\bf Theorem:} Let $\mathcal{B}=\{b_1, \dots, b_n\}$ and $\mathcal{C}=\{c_1, \dots, c_n\}$ be bases of a vector space $V$. Then there exists a unique $n\times n$ matrix $P$ such that $[x]_{\mathcal{C}}=P [x]_{\mathcal{B}}$.
\item We call $P$ the {\bf change of coordinates matrix from $\mathcal{B}$ to $\mathcal{C}$}. In general $$P= [[b_1]_\mathcal{C} \; \dots \; [b_n]_\mathcal{C}].$$ Also $P^{-1}$ is the change of coordinates matrix from $\mathcal{C}$ to $\mathcal{B}$.

\end{itemize}



\newpage

\begin{enumerate}



\Problem Let $V$ be a vector space with basis $\mathcal{B}=\{b_1, \dots, b_n\}$ and $W$ be a vector space with basis $\mathcal{C}=\{c_1, \dots, c_m\}$. Given a linear transformation $T: V\to W$, we naturally seek an associated matrix to represent $T$ using the bases; namely we want an $m\times n$ matrix $A$ such that $[T(x)]_\mathcal{C}=A [x]_\mathcal{B}$ for all $x \in V$. How can we find the matrix A?

\newpage

We call $A$ the {\bf matrix of $T$ relative to bases $\mathcal{B}$ and $\mathcal{C}$}. In general $A= [[T(b_1)]_\mathcal{C} \; \dots \; [T(b_n)]_\mathcal{C}]$.

\Problem Let $\mathcal{B}=\{b_1, b_2, b_3\}$ and $\mathcal{D}=\{d_1, d_2\}$ be bases for $V$ and $W$ respectively. Let $T: V \to W$ be the linear transformation such that $T(b_1)=3d_1-5d_2$, $T(b_2)=-d_1+6d_2$, and $T(b_3)=4d_2$. Find the matrix of $T$ relative to $\mathcal{B}$ and $\mathcal{D}$.

\vfill

\Problem Let $T: \mathbb{P}_2 \to \mathbb{P}_2$ by $T(p)=\frac{dp}{dt}$. Find the matrix of $T$ relative to the standard basis.

\vfill

\newpage


\Problem Let $A=\begin{bmatrix} 7 & 2\\ -4 & 1\\ \end{bmatrix}$. Let $T: \mathbb{R}^2 \to \mathbb{R}^2$ given by $T(x)=Ax$ and $\mathcal{B}=\left \{\begin{bmatrix} -1\\ 1\\ \end{bmatrix}, \begin{bmatrix} -1\\ 2\\ \end{bmatrix} \right\}.$
\Part Find the matrix of $T$ relative to $\mathcal{B}$.

\vfill

\Part Find $A^{2019}$.

\vfill 

{\bf Diagonal Representation Theorem:} Suppose $A= P D P^{-1}$ where $D$ is a diagonal matrix. If $\mathcal{B}$ is the basis for $\mathbb{R}^n$ formed by the columns of $P$, then $D$ is the matrix of $T$ relative to $\mathcal{B}$. 

\Problem Prove the theorem.

\vfill

\begin{defn}
If $A$ and $B$ are $n\times n$ matrices, then $A$ is {\bf similar} to $B$ if there exists an invertible matrix $P$ such that $A=PB P^{-1}$.
\end{defn}

\begin{defn}
If $A$ is an $n \times n$ matrix, then $\lambda \in \mathbb{R}$ is an {\bf eigenvalue} of $A$ if there exists a non-zero vector $v$, called an {\bf eigenvector} corresponding to $\lambda$, such that $$Av=\lambda v.$$
\end{defn}

\newpage

%\newpage
%
%\begin{center}
%{\bf Review Questions}
%\end{center}
%
%\Problem {True or False Questions}
%
%\begin{enumerate}
%\item  If $A$ and $B$ are $n$ by $n$ matrices with $\det(A)=2$ and $\det(B)=3$, then $\det(A+B)=5$.
%\item $k \text{det}(A) = \text{det}(kA)$
%%\item $AB$ is invertible if and only if $A$ and $B$ are invertible.
%\item If $A$ is row-equivalent to $B$, then \text{det}(A)=\text{det}(B).
%\item If $A$ is invertible and $AB=0$, then $B=0$.
%%\item If $T^2(x)=x$, then $T(x)=x$.
%%\item $S$ and $T$ linear transformations.  If $S \circ T$ is one-to-one, then T is one-to-one.  Must $S$ be one-to-one?
%%\item $S$ and $T$ linear transformations.  If $S \circ T$ is onto, then S onto.  Must $T$ be onto?
%\item If $A^2=I$, then $A=I$ or $A=-I$.
%\item If $A^2=0$, then $A=0$
%%\item there exists $u,w,u \in \mathbb{R}^3$ linearly dependent, so that none of the 3 is a multiple of another.
%\item $\text{det}(A) = \text{det}(B)$ implies $A-B$ is not invertible.
%%\item Let $M$ be an $n$ by $n$ matrix with $n$ odd.  If $M^{T}=-M$, then $M$ is not invertible.
%%\item If there exists a matrix $C$ so that $A= C^{-1} B C$, then \text{det}(A)=\text{det}(B).
%%\item Let $T : \mathbb{R}^3 \to \mathbb{R}^3$ be rotation around $z$-axis, and $A$ be the standard matrix for $T$, then $\det{A}=1$.
%%\item Let $T : \mathbb{R}^3 \to \mathbb{R}^3$ be projection onto the $yz$-plane, and $A$ be the standard matrix for $T$, then $\det{A}=0$.
%%\item If $A$ is invertible and $A^{T}=-A$, then A is an $n$ by $n$ matrix where $n$ is even.
%\item Let $T(x) =Ax$ and $S(x)=Bx$ and $A \sim B$, then $\text{im}(T) = \text{im}(S)$.
%\item Suppose $A$ and $B$ have the same RREF, then the span of the columns of $A$ equals the span of the columns of $B$.
%%\item If $Ax=b$ has a unique solution, then $A$ is a square matrix.
%%\item If $A^2+4A+3I=0$, and $A$ is 3 by 3, then $A$ is invertible.
%%\item If $A$ is $n$ by $n$ so that $A^2-A=0$, then $A=I$ or $A=0$.
%%\item There is a 2 by 2 matrix $A$ with real entries so that $A^2 = -I$.
%%\item There exists invertible 2 by 2 matrices $A$ and $B$ so that $\det(A+B)=\det(A)+\det(B)$.
%%\item  For any 2 by 2 matrix A, there exists $\lambda \in \mathbb{R}$ so that $A+ \lambda I$ is singular.
%\item  There is a linear transformation $T : \mathbb{R}^2 \to \mathbb{R}^2$ so that $\{ x : T(x)=0 \} = \{ T(x) : x \in \mathbb{R}^2 \}.$
%\item   If $u, v, w$ are pairwise linearly independent, then $u, v, w$ are linearly independent. 
%\end{enumerate}
%
%\Problem 
%\begin{enumerate}	
%\item		Let $T\colon \R^3 \rightarrow \R^3$ be the transformation that scales all vectors by a factor of 3.
%Let $A$ be the standard matrix for $T$.  What is $\det(A)$?
%%\begin{solution}
%%	We have that $T(\vec e_{1})=(3,0,0)$, $T(\vec e_{2})=(0,3,0)$ and $T(\vec e_{3})=(0,0,3)$.
%%	So the standard matrix for $T$ is $A=\begin{bmatrix}3&0&0\\0&3&0\\0&0&3\end{bmatrix}$.
%%	Since $A$ is triangular, the determinant of $A$ is the product of the diagonal entries: $\det(A)=27$.
%%\end{solution}
%%\blank
%%\part[3]
%\item Let $T\colon \R^3 \rightarrow \R^3$ be the transformation that reflects across the plane $x=z$.
%Let $A$ be the standard matrix for $T$.  What is $\det(A)$?
%%\begin{solution}
%%	We have that $T(\vec e_{1})=\vec e_3, T(\vec e_{2})=\vec e_{2}$ and $T(\vec e_{3})=\vec e_{1}$.
%%	So the standard matrix for $T$ is
%%	$A=\begin{bmatrix}0&0&1\\0&1&0\\1&0&0\end{bmatrix}.$ Since $A$ is obtained from the identity
%%	matrix by exchanging the first and third columns, $\det(A)=-1$.
%%\end{solution}
%%\blank
%%\part[3]
%\item Let $T\colon \R^3 \rightarrow \R^3$ be the transformation that projects to the plane $x=z$.
%Let $A$ be the standard matrix for $T$. What is $\det(A)$?
%%\begin{solution}
%%	Since $T$ is not onto $\R^{3}$, $T$ is not invertible. So the standard matrix $A$ for $T$ is not
%%	invertible. Thus, $\det(A)=0$.
%%\end{solution}
%%\blank
%\end{enumerate}
%
%\Problem Let \begin{equation*}
%A=\begin{bmatrix}
%1 & 0 & 1\\
%3 & 3 & 4\\
%2 & 2 & 3\\
%\end{bmatrix}.
%\end{equation*}  Find $A^{-1}$.
%
%\Problem 
%\begin{enumerate}
%	\item Prove that if $T$ is one-to-one and $\vec{u_1}, \cdots, \vec{u_p}\in \mathbb{R}^n$ are linearly independent, then $T(\vec{u_1}), \cdots, T(\vec{u_p})$ are linearly independent in $\mathbb{R}^m$.
%	\item In part (a), is the hypothesis that $T$ is one-to-one necessary?  Why or why not?
%\end{enumerate}
%
%\Problem Compute the determinant of the following matrices:
%\begin{enumerate}
%	\item
%	\begin{equation*}
%	A_1=\begin{bmatrix}
%	1 & 1\\
%	a_1 & 1\\
%	\end{bmatrix}.
%	\end{equation*}
%	\item 
%	\begin{equation*}
%	A_2=\begin{bmatrix}
%	1 & 1 & 1\\
%	a_1 & 1 & 1\\
%	a_2 & a_1 & 1\\
%	\end{bmatrix}.
%	\end{equation*}
%	\item Find the determinant of the following matrix.  Use induction to prove your formula is correct.
%	\begin{equation*}
%	A_n=\begin{bmatrix}
%	1 & 1 & \cdots & 1 & 1\\
%	a_1 & 1 & \cdots & 1& 1\\
%	\vdots &  & &&\vdots\\
%	a_{n-1} & a_{n-2} &  \ddots &  1    & 1\\
%	a_{n}     & a_{n-1} & \cdots & a_1 & 1\\
%	\end{bmatrix}.
%	\end{equation*}
%\end{enumerate}
%\Problem 
%\begin{enumerate}
%\item Prove $|\mathbb{N}|=|\{0\} \cup \mathbb{N}|$.
%\item Prove $|\mathbb{N}|=|\{-2,-1\} \cup \mathbb{N}|$.
%\item Prove $|[a,b]|=|[c,d]|$.
%\end{enumerate}
%
%\Problem Let  $T \colon \mathbb{P}_3 \rightarrow M_{2\times 2}$ be given by
%			$$ T(a_0 + a_1 t + a_2 t^2 + a_3 t^3) = \begin{bmatrix} a_0-a_1 & a_2 - a_3 \\ a_1-a_0 & a_3 - a_2 \end{bmatrix}. $$ Find bases for the kernel and range of $L$.  Justify your answer.  Find the standard matrix of $L$ relative to the standard bases.
%\Problem Let $U_{2 \times 2}$ be the space of upper triangular 2 by 2 matrices. Let $T : \mathbb{P}_2 \to U_{2 \times 2}$ be a linear transformation given by $$T(f) = \begin{bmatrix} f(0) & f'(1)\\  0 & f''(2)\\ \end{bmatrix}.$$   Find the matrix of $T$ relative to the standard bases.  Find the eigenvalues and eigenspaces of the matrix of $T$ relative to the standard bases.
%
%
%\Problem Let $T$ be the linear transformation $\mathbb R^3 \to \mathbb R^3$ with matrix
%	\[
%	A = \begin{bmatrix}
%	0 &
%	1 &
%	0 \\
%	
%	0 &
%	1 &
%	0 \\
%	
%	-1 &
%	1 &
%	1
%	\end{bmatrix}.
%	\]
%	
%	\begin{enumerate}
%		\item Compute the dimension of the kernel $T$ and a basis for the kernel.
%		
%		\item Compute the rank of $A$ and a basis for the image of $T$.
%	\end{enumerate}
%	
%%\Problem
%%	Let
%%	$$
%%	A = \begin{bmatrix}-10 &-18\\9&17\end{bmatrix}.
%%	$$
%%	\begin{enumerate}
%%		\item Compute the eigenvalues of $A$.
%%		\item Compute a basis for each eigenspace of $A$.
%%		\item Compute a matrix $B$ such that $B^3=A$.
%%	\end{enumerate}
%	
%\Problem Let $T:\mathbb{P}_2\to\mathbb{R}^3$ be defined by $$T(p(t)) = \begin{bmatrix}p(1)\\p(2)\\p(3)\end{bmatrix}$$ for every polynomial $p(t) = a_2t^2 + a_1t + a_0$ in $\mathbb{P}_2$.
%	
%	\begin{enumerate}
%		\item Show that $T$ is a linear transformation.
%		\item Find a matrix $M$ such that $T(p(t)) = M[p(t)]_{\mathcal{B}}$ where $\mathcal{B} = \{1,t,t^2\}$.
%		\item Show that for all $c_1, c_2, c_3 \in \mathbb{R}$ there is a polynomial $p(t) \in \mathbb{P}_2$ such that $p(1) = c_1$, $p(2) = c_2$, and $p(3) = c_3$.
%	\end{enumerate}
%	
%
%	
%
%%\Problem
%%	\begin{enumerate}
%%		\item
%%		Give an example of a real $n \times n$ matrix $A$ with eigenvalues $\frac{-1+i\sqrt{3}}{2}$ and $\frac{-1-i\sqrt{3}}{2}$.
%%		\item
%%		Give a complete list of all diagonalizable matrices with characteristic polynomial $(\lambda - 2)^4$. You must prove your list is complete.
%%		\item
%%		Find the algebraic multiplicities, and the dimension of the eigenspaces, of all the eigenvalues of the linear transformation $D:\mathbb{P}_2\to\mathbb{P}_2$ which sends each polynomial $p(t) = a_2t^2 + a_1t+ a_0$ in $\mathbb{P}_2$ to its derivative $p'(t) = 2a_2t+a_1$.
%%		
%%	\end{enumerate}
%	
%%\Problem
%%	Let $\mathcal{B} = \{\mathbf{b}_1,\mathbf{b}_2,\mathbf{b}_3\}$ be a basis for $\mathbb{R}^3$. Suppose that $A$ and $B$ are $3 \times 3$ matrices such that $\mathbf{b}_1,\mathbf{b}_2,\mathbf{b}_3$ are eigenvectors of both $A$ and $B$. Show that $AB = BA$.
%	
%\Problem
%	Let $A$ be an $n\times n$ matrix which is not invertible, and define $K_r$ to be the nullspace of $A^r$.  That is, $K_1 = \mbox{Nul}(A)$, $K_2 = \mbox{Nul}(A^2)$, etc.
%	
%	\begin{enumerate}
%		\item Prove that, for each $j=1,2,3,\ldots$, $K_j$ is a subspace of $K_{j+1}$, so that 
%		$$
%		K_1\subseteq K_2 \subseteq K_3 \subseteq K_4 \subseteq\cdots.
%		$$
%		\item Prove that, if for some $r$ we have $K_r = K_{r+1}$, then $K_r=K_n$ for all $n\geq r$.
%	\end{enumerate}
%	\Problem More True or False Questions.
%	\begin{enumerate}
%		 \item There is a $2 \times 4$ matrix $A$ with rank $1$ and $$\mathrm{Nul}(A) = \mathrm{Span}\left\{\begin{bmatrix}1\\0\\0\\1\end{bmatrix},\begin{bmatrix}1\\0\\1\\0\end{bmatrix},\begin{bmatrix}1\\1\\0\\0\end{bmatrix},\begin{bmatrix}3\\1\\1\\1\end{bmatrix}\right\}.$$
%		\item Let $V$ and $W$ be vector spaces and $T : V \to W$ be a one-to-one linear transformation.  If $\{ T(v_1), \cdots, T(v_k)\}$ is a basis for $W$, then $\{ v_1, \cdots, v_k \}$ is a basis for $V$.
%		\item The set $\{A\in M_{n \times n}(\mathbb{R}) : A \text{ is invertible. } \}$ is a subspace of $M_{n \times n} (\mathbb{R})$.
%		\item Let $A$ be an $n$ by $n$ matrix such that the sum of each row is equal to 1.  Then $A-I$ is not invertible.
%%		\item Let $E_1$ and $E_2$ be eigenspaces of distinct eigenvlaues of a matrix $A$.  Then $E_1 \cap E_2 = \{ 0 \}$.
%		\item  Let $W_1$ and $W_2$ be 3-dimensional subspaces of $\mathbb{R}^5$, then $W_1$ and $W_2$ have a non-zero vector in common.
%		\item $V$, $W$ subspaces of $\mathbb{R}^n$, then $V \cup W$ is a subspace.
%		\item  Let $W_1$ and $W_2$ be subspaces of $\mathbb{R}^n$ such that $\dim(W_1)\leq \dim(W_2)$.  Then $W_1\subset W_2$.
%	%	\item If $x$ is an eigenvector of a matrix $A$, then $x$ is also an eigenvector of $A^3$.
%		\item If $A$ is an $n \times n$ matrix and $x, y$ are two eigenvectors of $A$, then $x + y$ is also an eigenvector of $A$.
%		\item If $V$ and $W$ are vector spaces of dimension $n$ and $T : V \to W$ is a one-to-one linear transformation, then $T$ is onto.
%%		\item If $A$ and $B$ both have the characteristic polynomial $-\lambda^3+7\lambda^2-10\lambda$, then $A$ and $B$ are similar matrices.
%	\end{enumerate}
%

\end{enumerate}

\begin{comment}

\newpage
\chead{\Large\textbf{Eigenvalues and Eigenvectors -- Answers / Solutions}}%
\lhead{}
\rhead{}
\thispagestyle{fancy}

\begin{enumerate}
  \Answer Recall the matrix $M$ of $T$ relative to bases $\mathcal{B}$ and $\mathcal{D}$ is given by $M=[[T(b_1)]_\mathcal{D} \; [T(b_2)]_\mathcal{D} \; [T(b_3)]_\mathcal{D} ]$. Since $T(b_1)= 3d_1-5d_2$, we get $[T(b_1)]_\mathcal{D}=\begin{bmatrix} 3\\ -5\\ \end{bmatrix}$. Since $T(b_2)= -d_1+6d_2$, we get $[T(b_2)]_\mathcal{D}=\begin{bmatrix} -1\\ 6\\ \end{bmatrix}$. Since $T(b_3)= 4d_2$, we get $[T(b_3)]_\mathcal{D}=\begin{bmatrix} 0\\ 3\\ \end{bmatrix}$. Thus the matrix of $T$ relative to bases $\mathcal{B}$ and $\mathcal{D}$ is given by $\begin{bmatrix} 3 & -1 & 0\\ -5  & 6 & 4\\ \end{bmatrix}.$
  
  \Answer Let $\mathcal{E}=\{1, t, t^2\}$ be the standard basis for $\mathbb{P}_2$. Recall the matrix of $T$ relative to $\mathcal{E}$ is given by $$[T]_\mathcal{E} =[ [T(1)]_\mathcal{E} \; T(t)]_\mathcal{E} \; T(t^2)]_\mathcal{E}]$$. We compute the columns of the matrix: $T(1)=0$, and hence $[T(1)]=\begin{bmatrix} 0 \\ 0\\ 0\\ \end{bmatrix}$; $T(t)=1$, and hence $[T(t)]=\begin{bmatrix} 1 \\ 0\\ 0\\ \end{bmatrix}$; and $T(t^2)=2t$, and hence $[T(t^2)]=\begin{bmatrix} 2 \\ 0\\ 0\\ \end{bmatrix}$. Thus $$[T]_\mathcal{E} = \begin{bmatrix}
  0 & 1 & 0\\
  0 & 0 & 2\\
  0 & 0 & 0\\
  \end{bmatrix}.$$
  
  \Answer 
  
  \begin{proof} Let $P=[b_1 \; \dots \; b_n]$, then $\mathcal{B}=\{b_1,\dots, b_n\}$ is a basis for $\mathbb{R}^n$ since $P$ is invertible. We now compute the matrix of $T$ relative to $\mathcal{B}$.
  \begin{align*}
  [T]_\mathcal{B} & = [ [T(b_1)]_\mathcal{B} \; \dots [T(b_n)]_\mathcal{B} ]\\
   &=  [ [Ab_1]_\mathcal{B} \; \dots [Ab_n]_\mathcal{B} ]\\
   &=  [  P^{-1} Ab_1 \; \dots P^{-1} Ab_n ]\\
   & = P^{-1} A [b_1 \; \dots \; b_n]\\
   &= P^{-1}AP\\
   &= P^{-1}(PDP^{-1})P=D.
  \end{align*}
  Thus $T$ has a diagonal representation if $A$ is similar to a diagonal matrix $D$.
  \end{proof}
  
  \Answer Let $A=\begin{bmatrix} 7 & 2\\ -4 & 1\\ \end{bmatrix}$. Let $T: \mathbb{R}^2 \to \mathbb{R}^2$ given by $T(x)=Ax$ and $\mathcal{B}=\left \{\begin{bmatrix} -1\\ 1\\ \end{bmatrix}, \begin{bmatrix} -1\\ 2\\ \end{bmatrix} \right\}.$
  \begin{enumerate}
 \item Recall the matrix of $T$ relative to basis $\mathcal{B}$ is given by $[T]_\mathcal{B}=[ [T(b_1)]_\mathcal{B} \; [T(b_2)]_\mathcal{B} ]$.
 
 First note $$T(b_1)=Ab_1=\begin{bmatrix} 7 & 2\\ -4 & 1\\ \end{bmatrix}\begin{bmatrix} -1\\ 1\\ \end{bmatrix}=\begin{bmatrix} -5\\ 5\\ \end{bmatrix}=5\begin{bmatrix} -1\\ 1\\ \end{bmatrix},$$ and hence $[T(b_1)]_\mathcal{B}=\begin{bmatrix} 5\\ 0\\ \end{bmatrix}$. 
 
  Second note $$T(b_2)=Ab_2=\begin{bmatrix} 7 & 2\\ -4 & 1\\ \end{bmatrix}\begin{bmatrix} -1\\ 2\\ \end{bmatrix}=\begin{bmatrix} -3\\ 6\\ \end{bmatrix}=3\begin{bmatrix} -1\\ 2\\ \end{bmatrix},$$ and hence $[T(b_2)]_\mathcal{B}=\begin{bmatrix} 0\\3\\ \end{bmatrix}$. Thus the matrix of $T$ relative to basis $\mathcal{B}$ is given by $$[T]_\mathcal{B}=\begin{bmatrix} 5 & 0\\ 0 & 3\\ \end{bmatrix}.$$
 \item By the previous part, we know that $A= P_\mathcal{B} [T]_\mathcal{B} P_\mathcal{B}^{-1}$. Thus 
 \begin{align*}
 A^{2019} & = A \cdot A \cdot  \cdots  \cdot A \text{     \quad \quad   \quad \quad     2019 times} \\
 & = (P_\mathcal{B} [T]_\mathcal{B} P_\mathcal{B}^{-1}) (P_\mathcal{B} [T]_\mathcal{B} P_\mathcal{B}^{-1}) \cdots (P_\mathcal{B} [T]_\mathcal{B} P_\mathcal{B}^{-1}) \text{     \quad \quad   \quad \quad     2019 times} \\
 & = P_\mathcal{B} [T]_\mathcal{B}  [T]_\mathcal{B} \cdots [T]_\mathcal{B} P_\mathcal{B}^{-1}\text{     \quad \quad   \quad \quad     2019 times} \\
 &= P_\mathcal{B} [T]_\mathcal{B}^{2019} P_\mathcal{B}^{-1}\\
 &= \begin{bmatrix} -1 & -1\\ 1 & 2\\ \end{bmatrix} \begin{bmatrix} 5^{2019} & 0\\ 0 & 3^{2019} \end{bmatrix} \begin{bmatrix} -1 & -1\\ 1 & 2\\ \end{bmatrix}^{-1}\\
 &= \begin{bmatrix} -3^{2019}+5^{2019} & \frac{-3^{2019}}{2}+5^{2019}\\ 2\cdot 3^{2019}-5^{2019} & 3^{2019}-5^{2019}\\ \end{bmatrix}.
 \end{align*}
 \end{enumerate}

\end{enumerate}  
\end{comment}

%\end{document}


%%% Local ables: 
%%% mode: latex
%%% TeX-master: t
%%% End: 
